\chapter{部品の袋}
% full version: chapters/01b_neurontypes.tex

\pencilfig{1}{\pencilart{1}}{部品の袋：ほとんどが青、少しの赤、そしてひと握りのその他。}

見た目がそっくりな部品が860億個入った袋を渡され、これでどんなコンピュータが組み立てられているのか調べてくれと頼まれたとしよう。どこから手をつけるだろうか。エンジニアなら部品の形をじっくり眺めたりはしない。こう問うはずだ。入力はいくつあるか、出力はいくつか、そして出力に何を出すのか。

この部品はニューロンと呼ばれ、驚くほど単純にできている。入力は数千本ある。ほかのニューロンからの配線がつながっているのだ。出力は1本しかない。ただしこの出力は木のように枝分かれし、同じ信号のコピーが数千の受け手に届く。千人の話を同時に聞き、ひとことだけ答え、その答えをまた千人が聞く人を想像してほしい。

その信号とは何か。短いカチッという音である。ニューロンはときどきスパイクを出す。長さ約1000分の1秒の電圧の急上昇だ。1個のニューロンのカチッはどれも同じで、同じ太鼓を打つ音のように区別がつかない。だからニューロンが伝えられることはすべて、カチッの大きさではなく、\emph{いつ}鳴るかに書き込まれている。頻繁か、まばらか、すぐか、遅れてか。

ニューロンの内部では、考えうる最も単純な計算が行われている。入ってくる信号を足し合わせ、合計が閾値を超えたら自分もカチッと鳴る。ある入力は合計を押し上げ、別の入力は押し下げる。

\section*{袋の仕分け方}

出力配線の先端は、隣のニューロンに届いても触れてはいない。両者の間には細いすき間が残り、そこを信号が化学的に渡る。配線が特別な物質、つまり神経伝達物質をひとつまみ放出し、隣のニューロンがそれを受け取るのだ。そして幸運なことに、ほとんどのニューロンは配線のどの先端からも同じ物質を放出する。だからニューロンは、\emph{何を}放出するかで仕分けられる。

ニューロンの種類は、その物質の英語名の最初の4文字で呼ぶことにする。グルタミン酸（glutamate）を放出するニューロンは\nt{GLUT}。$\gamma$-アミノ酪酸（GABA）を放出するものは\nt{GABA}。以下、\nt{DOPA}（ドーパミン）、\nt{SERO}（セロトニン）、\nt{NORA}（ノルアドレナリン）、\nt{ACET}（アセチルコリン）などが続く。

袋を仕分けてみると、ひどく偏った構成になっている。1000個のニューロンのうち約960個が\nt{GLUT}だ。その信号は受け手をカチッの方へ押す。これが「プラス」である。約36個が\nt{GABA}で、その信号は受け手を閾値から遠ざける。これが「マイナス」だ。大脳皮質ではその割合が高く、5分の1から4分の1になる。残りの種類はすべて合わせても、10万個に数個。大海の一滴である。

\section*{電話と放送}

ところがこの一滴は、まったく違う仕組みで働いている。\nt{GLUT}や\nt{GABA}ニューロンの配線は、電話回線のように決まった相手につながっていて、信号を受け取るのは呼び出された相手だけだ。一方、\nt{DOPA}、\nt{SERO}、\nt{NORA}、\nt{ACET}の小さな集団は放送局のように働く。その配線は脳の広大な領域に広がり、物質はすき間ではなく周囲の空間に流れ出ることも多く、同じメッセージを数百万の細胞が同時に聞く。

電話回線を流れるのはデータだ。画像のここに線がある、この高さの音が鳴っている、といった情報である。放送で流れるのは全体の設定だ。全体の音量はどれくらいか、今は学習すべきか、もう眠る時間か。どのコンピュータも同じで、メモリセルは何十億もあるが、制御レジスタはひと握りしかない。それでも、それなしではマシンは動かない。

\section*{意味を決めるのは受け手}

ひとつ微妙な点がある。物質そのものは、自分が「プラス」なのか「マイナス」なのかを知らない。分子は鍵である。鍵が差し込まれたとき何が起きるかは錠、つまり受け手の側にある特別なタンパク質が決める。同じドーパミンが、ある受け手は励まし、別の受け手は抑える。どんな錠がついているかしだいだ。同じ放送を、聞き手によって違う意味に受け取るようなものである。

\section*{「予想よりよかった」という信号}

最も有名な放送チャネルはドーパミンだ。それは何を伝えているのか。こんな実験がある。動物に不意にジュースを一滴与えると、ドーパミンニューロンはカチッの連打で応える。次に、ジュースの前に毎回電球をつける。動物は電球がジュースの予告だと覚え、今度は連打が電球に対して起こり、ジュースそのものには起こらなくなる。そして電球のあとにジュースを与えないと、ドーパミンニューロンは期待していた瞬間に\emph{黙り込む}。

\pencilfig{101}{%
% three rows: a spike raster of a DOPA neuron; lamp (amber) and juice (blue drop) above the axis
\foreach \row/\y in {1/0,2/-1.05,3/-2.1}{
  \draw[pen] (0,\y) -- (6.2,\y);
  \foreach \s in {0.3,0.75,1.1,2.15,2.6,3.05,3.45,3.9,4.45,4.95,5.4,5.85}{
    \pgfmathtruncatemacro{\gap}{(\row==3)&&(\s>3.6)&&(\s<4.7)}
    \ifnum\gap=0 \draw[pcGray, line width=0.7pt] (\s,\y) -- (\s,\y+0.3);\fi}}
% row 1: juice without a lamp -> burst at the juice
\foreach \s in {4.0,4.08,4.16,4.24,4.33}{\draw[pcAmber, line width=0.8pt] (\s,0) -- (\s,0.45);}
\draw[dotpen=pcBlue, wash=pcBlue] (4.15,0.72) circle (0.09); \draw[pcBlue, line width=0.6pt] (4.07,0.76) -- (4.15,0.92) -- (4.23,0.76);
% row 2: lamp -> burst at the lamp, nothing at the promised juice
\foreach \y in {-1.05,-2.1}{
  \foreach \s in {1.5,1.58,1.66,1.74,1.83}{\draw[pcAmber, line width=0.8pt] (\s,\y) -- (\s,\y+0.45);}
  \draw[dotpen=pcAmber, wash=pcAmber] (1.65,\y+0.72) circle (0.1);
  \foreach \a in {30,90,150}{\draw[pcAmber, line width=0.5pt] ($(1.65,\y+0.72)+(\a:0.14)$) -- ($(1.65,\y+0.72)+(\a:0.24)$);}}
\draw[dotpen=pcBlue, wash=pcBlue] (4.15,-0.33) circle (0.09); \draw[pcBlue, line width=0.6pt] (4.07,-0.29) -- (4.15,-0.13) -- (4.23,-0.29);
% row 3: lamp, no juice -> a gap where the juice was expected
\draw[pcBlue, line width=0.6pt, densely dotted] (4.15,-1.38) circle (0.09); \draw[pcBlue, line width=0.6pt, densely dotted] (4.07,-1.34) -- (4.15,-1.18) -- (4.23,-1.34);
\draw[penlite=pcRed] (3.65,-2.22) -- (3.65,-2.28) -- (4.65,-2.28) -- (4.65,-2.22);
\node[lbl, text=pcRed, anchor=north] at (4.15,-2.3) {途切れ};
\node[lbl, anchor=east, align=right] at (-0.1,0.15) {電球なしで\\ジュース};
\node[lbl, anchor=east, align=right] at (-0.1,-0.9) {電球、\\次にジュース};
\node[lbl, anchor=east, align=right] at (-0.1,-1.95) {電球、\\ジュースなし};
}{予想外のジュースへの連打。ジュースを予告する電球への連打。そして、約束のジュースが来なかった瞬間の途切れ。}

三つの場合すべてを説明する規則はこうだ。ドーパミンが伝えているのは「よい」ではなく「予想よりよい」である。予想外のジュースはうれしい驚き。予告されたジュースは驚きなし。予告されたのに来ないジュースは、いやな驚き。自分の誤りから学ぶプログラムが使うのも、まさにこの量である。これについては後でまた触れる。

\section*{持ち帰ってほしいこと}

脳は、見た目がそっくりな部品の巨大なネットワークである。そのほとんどは、データを運ぶ電話網の「プラス」と「マイナス」だ。そしてごくわずかな部品が、ネットワーク全体をまとめて調整する放送局である。次の章では、「プラス」だけで何が作れるかを確かめる。

\fullversion{1}
